% ch42_derived_functors.tex — Volume IV Chapter 6

\chapter{Derived Functors and Quillen Equivalences}

\begin{overviewbox}
A functor between model categories rarely descends to homotopy categories
directly: a typical $F : \mathcal{C} \to \mathcal{D}$ fails to preserve
weak equivalences. The standard fix — \term{derive} $F$ — is to precompose
with cofibrant or fibrant replacement before evaluating. The resulting
\term{total left} or \term{total right derived functor} $\mathbb{L}F,
\mathbb{R}F : \mathrm{Ho}(\mathcal{C}) \to \mathrm{Ho}(\mathcal{D})$ is
the homotopy-coherent extension.

A \term{Quillen adjunction} $F \dashv G$ is a pair compatible with the
model structures: $F$ preserves cofibrations and trivial cofibrations,
$G$ preserves fibrations and trivial fibrations. Such adjunctions induce
\term{derived adjunctions} $\mathbb{L}F \dashv \mathbb{R}G$ on homotopy
categories. A Quillen adjunction is a \term{Quillen equivalence} when the
derived adjunction is an equivalence — the strongest comparison between
two model categories.

This chapter develops the derived-functor machinery and applies it to
prove the realization–singular Quillen equivalence (Theorem~41.\ref{thm:realization-sing-quillen}).
The chapter closes with two previews of Volume~IV's later applications:
the bar construction as a derived left Kan extension along $\Delta_+
\to \mathcal{V}$ (foreshadowing $\infty$-monad bar resolutions of
Chapter~51), and derived tensor / derived hom as the entry points to
homological algebra (Chapter~54). The phase-one foundational machinery
ends here; Chapter~43 turns to quasi-categories as the unifying
$\infty$-categorical model.
\end{overviewbox}


% ═══════════════════════════════════════════════════════════════════════════
\section{Total Derived Functors}
\label{sec:total-derived}
% ═══════════════════════════════════════════════════════════════════════════

\subsection{Informally}

Most functors do not respect weak equivalences. Tensoring a chain complex
with a non-flat module, for instance, does not preserve quasi-isomorphisms:
if $C \xrightarrow{\sim} C'$ is a quasi-isomorphism and $M$ is not flat,
$C \otimes M$ and $C' \otimes M$ may have different homology. The
\term{derived tensor product} $- \otimes^{\mathbb{L}} M$ fixes this by
first replacing the input chain complex by a quasi-isomorphic flat
complex (a \emph{cofibrant replacement}) and then tensoring.

The pattern is universal: to derive a functor that fails to respect
weak equivalences, replace inputs by cofibrant (or outputs by fibrant)
objects before evaluating. The result depends only on the weak-equivalence
class of the input, hence descends to $\mathrm{Ho}(\mathcal{C})$.

\subsection{Formalizing}

\begin{definition}[Total left and right derived functor]
\label{def:total-derived}
Let $\mathcal{C}, \mathcal{D}$ be model categories and $F : \mathcal{C}
\to \mathcal{D}$ a functor.

The \term{total left derived functor} of $F$, when it exists, is the
left Kan extension of the composite $\mathcal{C} \xrightarrow{F}
\mathcal{D} \xrightarrow{\gamma_{\mathcal{D}}} \mathrm{Ho}(\mathcal{D})$
along the localization $\gamma_{\mathcal{C}} : \mathcal{C} \to
\mathrm{Ho}(\mathcal{C})$:
\begin{equation*}
  \mathbb{L} F \;=\; \mathrm{Lan}_{\gamma_{\mathcal{C}}} (\gamma_{\mathcal{D}} F).
\end{equation*}

The \term{total right derived functor} $\mathbb{R} F$ is dually the right
Kan extension.
\end{definition}

\begin{howtosay}
$\mathbb{L} F$ \sayit{``the total left derived functor of $F$''};
$\mathbb{R} F$ \sayit{``the total right derived functor of $F$''}. The
blackboard-bold notation distinguishes the homotopical version from $F$
itself. Some texts write $LF$ and $RF$, or $F^{\mathbb{L}}$ and $F^{\mathbb{R}}$.
\end{howtosay}

\begin{theorem}[Pointwise formulas via replacement]
\label{thm:derived-pointwise}
If $F$ sends weak equivalences between cofibrant objects to weak
equivalences in $\mathcal{D}$, then $\mathbb{L} F$ exists and
\begin{equation*}
  \mathbb{L} F (X) \;\cong\; F(Q X) \quad \text{in } \mathrm{Ho}(\mathcal{D}),
\end{equation*}
where $Q X$ is a cofibrant replacement (Theorem~41.\ref{thm:replacement}).
Dually, if $F$ sends weak equivalences between fibrant objects to weak
equivalences,
\begin{equation*}
  \mathbb{R} F (X) \;\cong\; F(R X).
\end{equation*}
\end{theorem}

\begin{proof}[Proof sketch]
For $\mathbb{L}F$: define $\mathbb{L}F (X) = F(Q X)$ on objects.
Well-definedness on $\mathrm{Ho}(\mathcal{C})$: if $X \xrightarrow{\sim} X'$
is a weak equivalence, the induced $Q X \xrightarrow{\sim} Q X'$ between
cofibrants becomes a weak equivalence under $F$ by hypothesis, hence is
an isomorphism in $\mathrm{Ho}(\mathcal{D})$. The left-Kan-extension
universality follows from the universal property of cofibrant replacement:
$Q X \xrightarrow{\sim} X$ is initial among weak-equivalences-out-of-cofibrants
into $X$. \qed
\end{proof}

\begin{lemma}[Ken Brown's lemma]
\label{lem:ken-brown}
If $F : \mathcal{C} \to \mathcal{D}$ sends trivial cofibrations between
cofibrant objects to weak equivalences, then $F$ sends \emph{all} weak
equivalences between cofibrant objects to weak equivalences — so
$\mathbb{L} F$ exists.
\end{lemma}

\begin{remark}[Ken Brown in practice]
Ken Brown's lemma reduces the hypothesis ``preserves weak equivalences
between cofibrants'' to ``preserves trivial cofibrations.'' That is exactly
the half of the Quillen-functor condition that left adjoints often
satisfy; the next section makes this systematic.
\end{remark}


% ═══════════════════════════════════════════════════════════════════════════
\section{Quillen Adjunctions}
\label{sec:quillen-adjunctions}
% ═══════════════════════════════════════════════════════════════════════════

\subsection{Formalizing}

\begin{definition}[Quillen adjunction]
\label{def:quillen-adjunction}
A \term{Quillen adjunction} between model categories $\mathcal{C}$ and
$\mathcal{D}$ is an adjoint pair $F : \mathcal{C} \rightleftarrows
\mathcal{D} : G$, $F \dashv G$, such that any (equivalently, all) of the
following hold:
\begin{enumerate}
  \item $F$ preserves cofibrations and trivial cofibrations.
  \item $G$ preserves fibrations and trivial fibrations.
  \item $F$ preserves cofibrations and $G$ preserves fibrations.
\end{enumerate}
We call $F$ \term{left Quillen} and $G$ \term{right Quillen}.
\end{definition}

\begin{theorem}[Equivalence of Quillen-adjunction conditions]
\label{thm:quillen-adj-equiv}
For an adjunction $F \dashv G$, the three conditions in
Definition~\ref{def:quillen-adjunction} are equivalent.
\end{theorem}

\begin{proof}[Proof sketch]
(1) $\Leftrightarrow$ (2): cofibrations are exactly $\mathrm{LLP}$ against
trivial fibrations; by the adjunction $F \dashv G$, a square involving
$F(i)$ and $p$ has a lift iff the adjoint square involving $i$ and $G(p)$
does. So $F$ preserves cofibrations iff $G$ preserves trivial fibrations.
Dually for trivial cofibrations and fibrations.

(1) $\Rightarrow$ (3) is trivial. (3) $\Rightarrow$ (1): if $F$
preserves cofibrations, then $F$ preserves LLP-against-trivial-fibrations;
combined with $G$ preserving fibrations, $F$ also preserves trivial
cofibrations (via $\mathrm{LLP}$-against-fibrations). \qed
\end{proof}

\begin{theorem}[Derived adjunction]
\label{thm:derived-adjunction}
A Quillen adjunction $F \dashv G$ induces an adjunction of total derived
functors
\begin{equation*}
  \mathbb{L} F \;\dashv\; \mathbb{R} G \quad : \quad \mathrm{Ho}(\mathcal{C}) \rightleftarrows \mathrm{Ho}(\mathcal{D}).
\end{equation*}
\end{theorem}

\begin{proof}[Proof sketch]
By Ken Brown's lemma, $F$ preserving trivial cofibrations between
cofibrants implies $\mathbb{L}F$ exists; dually for $\mathbb{R}G$. The
adjunction descends because the hom-set bijection $\mathcal{D}(FX, Y)
\cong \mathcal{C}(X, GY)$ is compatible with the equivalence relation
defining $\mathrm{Ho}$: replacing $X$ by $QX$ on the left and $Y$ by $RY$
on the right gives
\begin{equation*}
  \mathrm{Ho}(\mathcal{D})(\mathbb{L} F (X), Y) \;\cong\; \mathrm{Ho}(\mathcal{C})(X, \mathbb{R} G (Y)). \qed
\end{equation*}
\end{proof}

\begin{example_thm}[Realization–singular is Quillen]
\label{ex:realization-quillen}
The adjunction $\lvert{-}\rvert \dashv \mathrm{Sing}$ between
$\mathbf{sSet}_{\mathrm{Q}}$ and $\mathbf{Top}_{\mathrm{Q}}$ is a
Quillen adjunction. Verification: $\lvert{-}\rvert$ preserves
monomorphisms (cofibrations of $\mathbf{sSet}$) since it preserves
colimits and monos in $\mathbf{Top}$ are pullback-stable. It also
preserves horn inclusions $\Lambda^n_k \hookrightarrow \Delta^n$
(trivial cofibrations of $\mathbf{sSet}$), since their realizations are
deformation retracts — hence trivial cofibrations in $\mathbf{Top}$.
\end{example_thm}


% ═══════════════════════════════════════════════════════════════════════════
\section{Quillen Equivalences}
\label{sec:quillen-equivalences}
% ═══════════════════════════════════════════════════════════════════════════

\subsection{Formalizing}

\begin{definition}[Quillen equivalence]
\label{def:quillen-eqv}
A Quillen adjunction $F \dashv G$ is a \term{Quillen equivalence} if any
(equivalently, all) of the following hold:
\begin{enumerate}
  \item The derived adjunction $\mathbb{L} F \dashv \mathbb{R} G$ is an
        equivalence of categories on $\mathrm{Ho}(\mathcal{C}) \simeq
        \mathrm{Ho}(\mathcal{D})$.
  \item For every cofibrant $X \in \mathcal{C}$ and fibrant $Y \in
        \mathcal{D}$, a map $F X \to Y$ is a weak equivalence in
        $\mathcal{D}$ iff its adjunct $X \to G Y$ is a weak equivalence
        in $\mathcal{C}$.
  \item For every cofibrant $X \in \mathcal{C}$, the composite
        $X \to G F X \xrightarrow{G(\text{fib repl})} G R F X$
        is a weak equivalence in $\mathcal{C}$; and dually for $Y$ fibrant.
\end{enumerate}
\end{definition}

\begin{theorem}[Realization–singular Quillen equivalence]
\label{thm:real-sing-qeq}
The adjunction $\lvert{-}\rvert \dashv \mathrm{Sing}$ between
$\mathbf{sSet}_{\mathrm{Q}}$ and $\mathbf{Top}_{\mathrm{Q}}$ is a Quillen
equivalence.
\end{theorem}

\begin{proof}[Proof sketch]
Use condition (2). For $X \in \mathbf{sSet}$ cofibrant (every simplicial
set is) and $Y \in \mathbf{Top}$ fibrant (every space is), the question
is whether $\lvert X\rvert \to Y$ is a weak homotopy equivalence iff
$X \to \mathrm{Sing}\, Y$ is.

Both directions reduce to the unit and counit:
\begin{itemize}
  \item Unit $\eta_X : X \to \mathrm{Sing}\,\lvert X\rvert$. By a theorem of
        Milnor, this is a weak equivalence in $\mathbf{sSet}$ for every $X$:
        $\mathrm{Sing}\,\lvert X\rvert$ and $X$ have the same simplicial
        homotopy groups (Chapter~38's $\pi_n$ definition, via inner-horn
        filling).
  \item Counit $\epsilon_Y : \lvert\mathrm{Sing}\, Y\rvert \to Y$. This is
        a weak homotopy equivalence for every space $Y$ (Quillen's
        theorem), classically equivalent to $Y$ admitting a CW
        approximation.
\end{itemize}
Both unit and counit being weak equivalences is condition (3); hence the
adjunction is a Quillen equivalence. \qed
\end{proof}

\begin{remark}[The homotopy hypothesis, formalized]
Theorem~\ref{thm:real-sing-qeq} is the precise statement of the
\emph{homotopy hypothesis} we previewed in Chapter~39: simplicial sets
and topological spaces present equivalent homotopy theories. From this
point onward in Volume~IV, we may freely move between $\mathbf{sSet}$
and $\mathbf{Top}$ at the level of homotopy categories.
\end{remark}


% ═══════════════════════════════════════════════════════════════════════════
\section{Examples and Bridges}
\label{sec:bridges}
% ═══════════════════════════════════════════════════════════════════════════

\subsection{The bar construction as a derived left Kan extension}

\begin{example_thm}[Derived bar]
\label{ex:derived-bar}
For a monoid $M$ in a strict monoidal model category $(\mathcal{V},
\otimes, I)$, the bar construction $B_\bullet M : \Delta^{\mathrm{op}}
\to \mathcal{V}$ (Chapter~37 Corollary~\ref{cor:bar-construction})
gives a simplicial object. Its \term{geometric realization}
$\lvert B_\bullet M\rvert$ — the coend $\int^{[n]} M^{\otimes(n+1)} \otimes
\lvert\Delta^n\rvert$ — computes the derived colimit
\begin{equation*}
  \mathbb{L}\mathrm{colim}_{\Delta^{\mathrm{op}}} B_\bullet M \;\simeq\; \lvert B_\bullet M\rvert.
\end{equation*}
This is the model-categorical incarnation of ``the universal property of
$\Delta_+$ derived to its homotopical version.'' The $\infty$-monad
machinery of Chapter~51 makes this systematic.
\end{example_thm}

\subsection{Derived tensor and Ext}

\begin{example_thm}[Derived tensor product]
\label{ex:derived-tensor}
Let $R$ be a ring, $M$ a right $R$-module, $N$ a left $R$-module.
Cofibrantly resolve $M$ in the projective model structure on
$\mathbf{Ch}(\mathrm{Mod}\text{-}R)$ to obtain a chain complex
$P_\bullet \xrightarrow{\sim} M$ of projective modules. Then
\begin{equation*}
  M \otimes^{\mathbb{L}}_R N \;=\; P_\bullet \otimes_R N,
\end{equation*}
the \term{derived tensor product}, computes the Tor groups $\mathrm{Tor}_n^R(M, N)$
as $H_n(P_\bullet \otimes_R N)$.
\end{example_thm}

\begin{example_thm}[Derived hom]
\label{ex:derived-hom}
Dually, $\mathbb{R}\mathrm{Hom}_R(M, N) = \mathrm{Hom}_R(P_\bullet, N)$
where $P_\bullet \xrightarrow{\sim} M$ is a projective resolution
(or $N \xrightarrow{\sim} I^\bullet$ an injective resolution), with
homology computing $\mathrm{Ext}^n_R(M, N) = H^n \mathbb{R}\mathrm{Hom}_R(M, N)$.
\end{example_thm}

\begin{remark}[Phase 1 closes here]
With derived functors and Quillen equivalences in place, the
simplicial / model-categorical foundations of Volume~IV are complete.
The next phase (Chapters~43--48) builds the theory of quasi-categories
on this foundation. The key takeaway: $\infty$-category theory will
generalize \emph{both} ordinary categories (via the nerve) \emph{and}
simplicial homotopy types (via Kan complexes), unifying the two pictures
under a single combinatorial framework.
\end{remark}


% ═══════════════════════════════════════════════════════════════════════════
\section{Exercises}
\label{sec:derived-exercises}
% ═══════════════════════════════════════════════════════════════════════════

\begin{exercisesbox}
\begin{qa}
\qaitem{Define the total left derived functor $\mathbb{L}F$.}{The left Kan extension of $\gamma_{\mathcal{D}} \circ F$ along the localization $\gamma_{\mathcal{C}}$; pointwise $\mathbb{L}F(X) \simeq F(QX)$ when $F$ preserves weak equivalences between cofibrants.}
\qaitem{What does $\mathbb{R}F$ compute?}{Total right derived functor: $\mathbb{R}F(X) \simeq F(RX)$ via fibrant replacement, when $F$ preserves weak equivalences between fibrants.}
\qaitem{State Ken Brown's lemma.}{If $F$ preserves trivial cofibrations between cofibrant objects, it preserves all weak equivalences between cofibrant objects.}
\qaitem{Why is Ken Brown's lemma useful?}{It reduces the existence of $\mathbb{L}F$ to checking $F$ preserves trivial cofibrations — a much simpler condition than ``preserves all weak equivalences.''}
\qaitem{Define a Quillen adjunction.}{An adjunction $F \dashv G$ with $F$ left Quillen (preserves cof and trivial cof) and $G$ right Quillen (preserves fib and trivial fib).}
\qaitem{Why are the three conditions in Definition~\ref{def:quillen-adjunction} equivalent?}{Cofibrations are LLP against trivial fibrations, and LLP/RLP are exchanged by the adjunction. So $F$ preserves cof iff $G$ preserves trivial fib, etc.}
\qaitem{State the derived adjunction theorem.}{A Quillen adjunction $F \dashv G$ induces $\mathbb{L}F \dashv \mathbb{R}G : \mathrm{Ho}(\mathcal{C}) \rightleftarrows \mathrm{Ho}(\mathcal{D})$.}
\qaitem{Define a Quillen equivalence.}{A Quillen adjunction whose derived adjunction is an equivalence of homotopy categories.}
\qaitem{Equivalent characterization in terms of unit/counit?}{$X \xrightarrow{\eta} GFX \xrightarrow{G(\sim)} GRFX$ is a weak equivalence for every cofibrant $X$, and the dual for fibrant $Y$.}
\qaitem{Why is $\lvert{-}\rvert \dashv \mathrm{Sing}$ a Quillen equivalence?}{Unit $X \to \mathrm{Sing}\,\lvert X\rvert$ is a weak equivalence (Milnor); counit $\lvert\mathrm{Sing}\, Y\rvert \to Y$ is a weak equivalence (Quillen / CW approximation).}
\qaitem{Why is $\lvert{-}\rvert$ left Quillen?}{It preserves monomorphisms (so cofibrations of $\mathbf{sSet}$) and sends horn inclusions to deformation retracts (so trivial cofibrations).}
\qaitem{Why is $\mathrm{Sing}$ right Quillen?}{By adjunction with $\lvert{-}\rvert$ being left Quillen — or directly: $\mathrm{Sing}(p)$ is a Kan fibration iff $p$ is a Serre fibration.}
\qaitem{What does $\mathbb{L}F$ compute when $F$ already preserves weak equivalences?}{Just $F$ itself, descended to $\mathrm{Ho}(\mathcal{C})$. The derivation is trivial.}
\qaitem{When is $F : \mathcal{C} \to \mathcal{D}$ \emph{homotopical}?}{When $F$ preserves \emph{all} weak equivalences. Such $F$ has $\mathbb{L}F = \mathbb{R}F = F_{\mathrm{Ho}}$.}
\qaitem{What is the derived tensor product?}{$M \otimes_R^{\mathbb{L}} N = P_\bullet \otimes_R N$, where $P_\bullet \xrightarrow{\sim} M$ is a projective resolution. $H_n$ recovers $\mathrm{Tor}^R_n(M, N)$.}
\qaitem{What is derived hom?}{$\mathbb{R}\mathrm{Hom}_R(M, N) = \mathrm{Hom}_R(P_\bullet, N)$ for $P_\bullet$ projective resolution of $M$; cohomology recovers $\mathrm{Ext}^n_R(M, N)$.}
\qaitem{Why does Tor not just equal tensor at the chain level?}{Because tensoring with a non-flat module doesn't preserve quasi-isomorphisms; deriving via cofibrant replacement (projective resolution) corrects this.}
\qaitem{What is the homotopy colimit?}{$\mathrm{hocolim} = \mathbb{L}\mathrm{colim}$. For a simplicial diagram, $\mathrm{hocolim}_\Delta = \lvert{-}\rvert$ (geometric realization of the simplicial replacement).}
\qaitem{The bar construction as a derived colimit?}{$\lvert B_\bullet M\rvert \simeq \mathbb{L}\mathrm{colim}_{\Delta^{\mathrm{op}}} B_\bullet M$ — the realization of the bar simplicial object is the derived colimit.}
\qaitem{What's a Bousfield localization?}{A new model structure on the same underlying category and cofibrations, with a strictly larger class of weak equivalences. Not developed here, but central in advanced applications.}
\qaitem{Does every Quillen adjunction yield a Quillen equivalence?}{No — only those that, in addition, witness that no homotopy-theoretic information is lost. The realization–singular adjunction does; many others don't.}
\qaitem{What is a homotopy hypothesis?}{The claim that $\infty$-groupoids and homotopy types coincide. Theorem~\ref{thm:real-sing-qeq} is its model-categorical version: $\mathbf{sSet}_{\mathrm{Q}} \simeq_{\mathrm{Q}} \mathbf{Top}_{\mathrm{Q}}$.}
\qaitem{What's the next phase of Volume~IV?}{Quasi-categories (Chapter~43) — combining the nerve perspective (Chapter~39) with the Joyal model structure (Chapter~41 preview) to give a single combinatorial model of $\infty$-categories.}
\end{qa}
\end{exercisesbox}


% ═══════════════════════════════════════════════════════════════════════════
\section*{Chapter Summary}
\addcontentsline{toc}{section}{Chapter Summary}
% ═══════════════════════════════════════════════════════════════════════════

\begin{summarybox}
\textbf{Total derived functor.}\quad $\mathbb{L}F = \mathrm{Lan}_\gamma
(\gamma_{\mathcal{D}} F)$ — pointwise $\mathbb{L}F(X) \simeq F(QX)$ via
cofibrant replacement. Dually $\mathbb{R}F(X) \simeq F(RX)$.

\textbf{Ken Brown.}\quad Preserving trivial cofibrations between
cofibrants suffices for $\mathbb{L}F$ to exist. The right-handed dual
gives $\mathbb{R}F$.

\textbf{Quillen adjunction.}\quad $F \dashv G$ with $F$ left Quillen
(preserves cof and trivial cof) and $G$ right Quillen (dually). The
conditions are exchanged by adjunction. Induces derived adjunction
$\mathbb{L}F \dashv \mathbb{R}G$ on homotopy categories.

\textbf{Quillen equivalence.}\quad Derived adjunction is an equivalence.
Realization–singular $\lvert{-}\rvert \dashv \mathrm{Sing}$ is the
canonical example: simplicial sets and topological spaces are
Quillen-equivalent. This is the homotopy hypothesis formalized.

\textbf{Bridges.}\quad The bar construction realizes as a derived
colimit; derived tensor computes Tor; derived hom computes Ext. Phase~1
foundations close here; Phase~2 (Chapter~43+) builds quasi-categories
on top.
\end{summarybox}


% ═══════════════════════════════════════════════════════════════════════════
\section*{Formal Restatement}
\addcontentsline{toc}{section}{Formal Restatement}
% ═══════════════════════════════════════════════════════════════════════════

\begin{formalrestatement}
\textbf{Total derived functors.}\quad
\begin{equation*}
  \mathbb{L} F \;=\; \mathrm{Lan}_{\gamma_{\mathcal{C}}} (\gamma_{\mathcal{D}} F), \qquad
  \mathbb{R} F \;=\; \mathrm{Ran}_{\gamma_{\mathcal{C}}} (\gamma_{\mathcal{D}} F).
\end{equation*}
Pointwise:
$\mathbb{L} F (X) \simeq F(QX)$, $\mathbb{R} F(X) \simeq F(RX)$.

\medskip
\textbf{Ken Brown's lemma.}\quad If $F$ sends trivial cofibrations between
cofibrants to weak equivalences, $F$ sends all weak equivalences between
cofibrants to weak equivalences.

\medskip
\textbf{Quillen adjunction.}\quad $F : \mathcal{C} \rightleftarrows
\mathcal{D} : G$ with $F$ preserving $\mathcal{C}\mathrm{of}$ and
$\mathcal{W} \cap \mathcal{C}\mathrm{of}$ (equivalently $G$ preserving
$\mathcal{F}\mathrm{ib}$ and $\mathcal{W} \cap \mathcal{F}\mathrm{ib}$).

\medskip
\textbf{Derived adjunction.}\quad $\mathbb{L} F \dashv \mathbb{R} G :
\mathrm{Ho}(\mathcal{C}) \rightleftarrows \mathrm{Ho}(\mathcal{D})$.

\medskip
\textbf{Quillen equivalence.}\quad Quillen adjunction whose derived
adjunction is an equivalence. Equivalently: for $X$ cofibrant, $Y$
fibrant, $FX \to Y$ is a w.e.\ in $\mathcal{D}$ iff $X \to GY$ is a w.e.\
in $\mathcal{C}$.

\medskip
\textbf{Homotopy hypothesis (Theorem).}\quad $\lvert{-}\rvert :
\mathbf{sSet}_{\mathrm{Q}} \rightleftarrows \mathbf{Top}_{\mathrm{Q}} :
\mathrm{Sing}$ is a Quillen equivalence; $\mathrm{Ho}(\mathbf{sSet})
\simeq \mathrm{Ho}(\mathbf{Top})$.

\medskip
\textbf{Derived algebra.}\quad $\mathrm{Tor}_n^R(M, N) = H_n(M \otimes_R^{\mathbb{L}} N)$;
$\mathrm{Ext}_R^n(M, N) = H^n \mathbb{R}\mathrm{Hom}_R(M, N)$. Cofibrant
replacement $=$ projective resolution; fibrant replacement on opposite
side $=$ injective resolution.
\end{formalrestatement}
